\begin{lemma}[Stability of $\gamma_n$ under an orthogonal mixing of two independent copies]
Fix $n \in \mathbb{N}$ and let $a, b \in \mathbb{R}$ satisfy
\[
  a \geq 0, \qquad b \geq 0, \qquad a^2 + b^2 = 1 .
\]
Let $\gamma_n$ be the Gaussian measure of \noderef{GaussMeasure} on
$\mathbb{R}^n$, and consider the map
\[
  \Phi_{a,b} : \mathbb{R}^n \times \mathbb{R}^n \to \mathbb{R}^n,
  \qquad \Phi_{a,b}(x,y) = a\,x + b\,y .
\]
Then the pushforward of the product measure $\gamma_n \otimes \gamma_n$ along
$\Phi_{a,b}$ is $\gamma_n$ again:
\[
  (\Phi_{a,b})_* \bigl( \gamma_n \otimes \gamma_n \bigr) = \gamma_n .
\]

Every measure here is Gaussian: the source is the product of two copies of
$\gamma_n$ on $\mathbb{R}^n \times \mathbb{R}^n$ and the target is $\gamma_n$
itself, with no Lebesgue measure appearing except inside the density that
\noderef{GaussMeasure} already carries.  In probabilistic language: if $X$ and $Y$
are independent with law $\gamma_n$ and $a^2 + b^2 = 1$, then $aX + bY$ again has
law $\gamma_n$.  This is the property that lets a $\gamma_n$-distributed variable
be split as a contraction of itself plus an independent Gaussian of the
complementary variance, which is what makes the Ornstein--Uhlenbeck evolution
expressible as an average against an independent copy; the explicit size of the
resulting kernel is never needed, only that such a representation exists.

The hypotheses $a \geq 0$ and $b \geq 0$ are what the proof below actually uses,
at the one point where $\|a\,t\|$ is replaced by $a\,\|t\|$; see the remark at the
end of the proof.
\end{lemma}
\begin{proof}
\emph{Step 1: it suffices to compare characteristic functions.}
The map $\Phi_{a,b}$ is measurable, being built from scalar multiplication and
addition.  By \noderef{GaussIsProb} the measure $\gamma_n$ is a probability
measure, hence so is the product $\gamma_n \otimes \gamma_n$, and hence so is its
pushforward along a measurable map.  Two probability measures on $\mathbb{R}^n$
with the same characteristic function are equal, so it is enough to prove
\begin{equation*}
  \widehat{(\Phi_{a,b})_*(\gamma_n \otimes \gamma_n)}(t)
    = \widehat{\gamma_n}(t)
  \qquad\text{for every } t \in \mathbb{R}^n . \tag{1}
\end{equation*}

\emph{Step 2: the left-hand side as an integral over the product.}
By the change-of-variables formula for a pushforward along a measurable map, the
characteristic function of the pushforward is the integral of the same exponential
composed with $\Phi_{a,b}$:
\[
  \widehat{(\Phi_{a,b})_*(\gamma_n \otimes \gamma_n)}(t)
    = \int_{\mathbb{R}^n \times \mathbb{R}^n}
        e^{\,i \langle a\,x + b\,y,\ t\rangle}
      \, d(\gamma_n \otimes \gamma_n)(x,y) .
\]

\emph{Step 3: the integrand factorizes.}
For all $x, y, t$, bilinearity of the inner product gives
\[
  \langle a\,x + b\,y,\ t\rangle
    = a\,\langle x, t\rangle + b\,\langle y, t\rangle
    = \langle x,\ a\,t\rangle + \langle y,\ b\,t\rangle ,
\]
so, since $e^{z+w} = e^z e^w$,
\begin{equation*}
  e^{\,i\langle a x + b y,\ t\rangle}
    = e^{\,i\langle x,\ a t\rangle} \cdot e^{\,i\langle y,\ b t\rangle} . \tag{2}
\end{equation*}
The two factors depend on $x$ alone and on $y$ alone respectively.

\emph{Step 4: Fubini for the product measure.}
Each factor in $(2)$ has modulus $1$, hence is integrable against the probability
measure $\gamma_n$ in its own variable.  The integral of a product of two such
one-variable functions against $\gamma_n \otimes \gamma_n$ is therefore the
product of the two integrals, and each of those is a characteristic function of
$\gamma_n$:
\[
  \int_{\mathbb{R}^n\times\mathbb{R}^n} e^{\,i\langle x, a t\rangle}
      e^{\,i\langle y, b t\rangle} \, d(\gamma_n\otimes\gamma_n)(x,y)
  = \widehat{\gamma_n}(a\,t)\ \widehat{\gamma_n}(b\,t) .
\]

\emph{Step 5: evaluate, and use $a^2 + b^2 = 1$.}
By \noderef{GaussCharFun}, $\widehat{\gamma_n}(u) = e^{-|u|^2/4}$ for every
$u \in \mathbb{R}^n$.  Now $|a\,t| = |a|\,|t| = a\,|t|$, the last equality because
$a \geq 0$, and likewise $|b\,t| = b\,|t|$ because $b \geq 0$.  Hence
\[
  \widehat{\gamma_n}(a t)\ \widehat{\gamma_n}(b t)
    = e^{-a^2|t|^2/4}\, e^{-b^2|t|^2/4}
    = e^{-(a^2+b^2)|t|^2/4}
    = e^{-|t|^2/4}
    = \widehat{\gamma_n}(t),
\]
using $e^{z}e^{w} = e^{z+w}$ for the middle step and $a^2 + b^2 = 1$ for the next.
This is $(1)$, and Step 1 then gives
$(\Phi_{a,b})_*(\gamma_n \otimes \gamma_n) = \gamma_n$.

Remark.  The conclusion in fact holds for arbitrary real $a, b$ with
$a^2 + b^2 = 1$, since Step 5 would then read $|a\,t| = |a|\,|t|$ and
$|a|^2 = a^2$ leaves the computation unchanged.  The sign hypotheses are stated
because they are what this proof route consumes --- they are used only to
rewrite $|a|$ as $a$ and $|b|$ as $b$ --- and not because the statement needs
them.
\end{proof}
